\documentclass[10pt]{article}

\usepackage[utf8]{inputenc}

\usepackage[T1]{fontenc}

\usepackage{amsmath}

\usepackage{amsfonts}

\usepackage{amssymb}

\usepackage[version=4]{mhchem}

\usepackage{stmaryrd}

\usepackage{bbold}



\begin{document}

имеющее О.П. $F_{\psi}$, то $F_{\psi}$ удовлетворяет уравнению Лиувилля



\[

\partial_{t} F_{\psi}+\left\{H, F_{\psi}\right\}=0

\]



$\left\{H, F_{\psi}\right\}-$ Пуассона скобка функций $H$ и $F_{\psi}$.



Лит.: [1] Карасе в М. В., Масло в В. П., «Успехи математич. $\begin{array}{ll}\text { наук», 1984, т. 39, в. 6, с. 115-73. } & \text { В. Е. Назайкинский. }\end{array}$



осцилляций сУБплотность - коэффициенты $G_{\psi}$ (обобщенная функция на фазовом пространстве $\mathbb{R}_{q, p}^{2 n}$ ) разложения по $h$ функции плотности Блохинцева-Вагнера $\rho_{\psi}^{h}$, отвечающей волновой функции $\psi$ (см. [1]):



\[

\rho_{\psi}^{h}=F_{\psi}+h G_{\psi}+O(h)

\]



где $F_{\psi}$ - осцилляиий плотность. Напр., О. с. для гауссова пакета



имеет вид:



\[

\psi(q)=(\pi h)^{-n / 4} \exp \left(-q^{2} / 2 h\right)

\]



\[

G_{\Psi}=\frac{1}{4} \Delta(\delta(q) \delta(p))

\]



где $\Delta$ - оператор Лапласа на $\mathbb{R}_{q, p}^{2 n}$.



Лит.: [1] Карасе в М. В., Масло в В. П., «Успехи математич. наук», 1984, т. 39, в. 6, с. 115-73. Ю. М. Воробвев. ОСцилляций ФРОНТ волновой Функции $\psi-$ подмножество фазового пространства, представляющее собой носитель соответствующей осцияляиий плотности $F_{\psi}$. Если $\psi$ - решение уравнения Шредингера



\[

-i h \partial_{t} \psi+H\left(x,-i h \partial_{x}\right) \psi=0,

\]



то О.ф. распространяется вдоль характеристик функции Гамильтона $H(x, p)$ (решений системы Гамильтона $\left.\dot{x}=H_{p}(x, p), \dot{p}=-H_{x}(x, p)\right)$.



См. также Осцилляций носитель. $\quad$ В. Е. Назайкинский. ОСцилляцИй ШлЕйФ - область, возникающая в результате распространения возмущений от точечных источников в средах с дисперсией. Типичными примерами являются среды, описываемые разностными, дифференциально-разностными или псевдодифференциальными уравнениями с малым параметром. Математич. описание этого явления на примере уравнения кристалла (см. Разностные уравнения; асимптотика решения)



\[

\begin{gathered}

h^{2} u_{t t}+L\left(x,-i h \partial / \partial_{x}\right) u=0, \\

L(x, p)=4 c^{2}(x) \sin ^{2} p / 2, x \in \mathbb{R}^{1},-\pi \leqslant p \leqslant \pi,

\end{gathered}

\]



$h$ - малый параметр, заключается в том, что областью осцилляций решения задачи Коши для этого уравнения с начальными условиями $\left.u\right|_{t=0}=h^{-1 / 2} \delta_{h}\left(x-x^{0}\right),\left.u_{t}\right|_{t=0}=0$ в момент времени $t$ является множество $\left\{x=X^{ \pm}\left(x^{0}, t, \omega\right)\right.$, $-\pi \leqslant \omega \leqslant \pi\}$, где $X^{ \pm}\left(x^{0}, t, \omega\right), P^{ \pm}\left(x^{0}, t, \omega\right)$ - решение гамильтоновой системы $\dot{x}=\partial H_{\omega}^{ \pm} / \partial p, \dot{p}=-\partial H_{\omega}^{ \pm} / \partial x$, удовлетворяющее начальным данным $\left.X^{ \pm}\right|_{t=0}=x^{0},\left.P^{ \pm}\right|_{t=0}=1$. Область осцилляций функции $\psi_{h}(x) \in L_{2}\left(\mathbb{R}^{n}\right)$ определяется как дополнение к множеству точек таких, что в их окрестности $\psi_{h} \in L_{2}$ и $\partial \psi_{h} / \partial x_{i} \in L_{2}, i=1, \ldots, n$, равномерно по $h$. В том случае, когда $\psi_{h}=K \varphi$, где $K-$ Маслова канонический оператор на лагранжевом многообразии $\Lambda \subset \mathbb{R}_{x}^{n} \times \mathbb{R}_{p}^{n}$, область осцилляций совпадает с проекцией осциляяций фронта на конфигурационное пространство $\mathbb{R}_{x}^{n}$. В указанном примере область осцилляций начального данного - точка $x^{0}$, а область осцилляций решения в момент времени $t$ - нек-рый отрезок прямой $\mathbb{R}^{1}$.



См. также Черенкова эффект в разностных уравнениях.



Лит.: [1] Масло в В. П., Операторные методы, М., 1973; [2] Маслов В. П., Дан илов В. Г., «Труды Математич. ин-та АН СССР», 1984, т. 166, с. 130-60; 1985, т. 167, с. 96-107; [3] До б рох о то в С.Ю., Же в андров П. Н., «Функциональный анализ и его приложения», 1985, т. 19, № 4, с. 43-54.



В. Г. Данилов, П. Н. Жевандров. ОТБОРА ПРАВИЛА - правила, определяющие возможные квантовые переходы для атомов, молекул, атомных ядер, взаимодействующу элементарных частиц и др. О. п. устанавливают, какие переходы разрешены (вероятность перехода велика) и какие запрещены строго (вероятность перехода равна нулю) или приближенно (вероятность перехода мала); соответствующие О. п. разделяют на строгие и приближенные. При характеристике состояний системы с помощью квантовых чисел О.п. определяют возможные изменения этих чисел при переходе рассматриваемого типа.



О. п. связаны с симметрией квантовых систем, то есть с неизменностью (инвариантностью) их свойств при определенных преобразованиях, в частности координат и времени, и с соответствующими сохранения законами. Переходы с нарушением строгих законов сохранения энергии, импульса (для реальных переходов), момента импульса, электрич. заряда и т. д. замкнутой системы абсолютно исключаются.



Для излучательных квантовых переходов между стационарными состояниями атомов и молекул очень важны строгие О. п. для квантовых чисел $J$ и $m_{J}$, определяющих возможные значения полного момента импульса $\boldsymbol{M}$ и его проекции $M_{z}\left[\boldsymbol{M}^{2}=\hbar^{2} J(J+1), M_{z}=\hbar m_{J}\right]$. Эти правила связаны с равноправием в пространстве всех направлений (для любой точки - сферическая симметрия) и всех направлений, перпендикулярных выделенной оси $z$ (а кс и ал ь на я с им ме т р и я), и соответствуют сохранению момента импульса и его проекции на ось $z$. Из законов сохранения полного момента импульса и его проекции для системы, состоящей из микрочастиц и из испускаемых, поглощаемых и рассеиваемых фотонов, следует, что при квантовом переходе $J$ и $m_{J}$ могут изменяться в случае электрич. и магнитных дипольных излучений лишь на 0 , \textbackslash pm 1 , а в случае электрич. квадрупольного излучения на 0 , $\pm 1, \pm 2$.



Другое важное О. п. связано с законом сохранения полной четности для изолированной квантовой системы (этот закон нарушается лишь слабым взаимодействием). Квантовые состояния атомов, всегда имеющих центр симметрии, а также тех молекул и кристаллов, к-рые имеют такой центр, делятся на четные и нечетные по отношению к пространственной инверсии (отражению в центре симметрии, то есть к преобразованию координат $x \rightarrow-x, y \rightarrow-y$, $z \rightarrow-z)$; в этих случаях справедлив так наз. альтернативный запрет для излучательных квантовых переходов: для электрич. дипольного излучения запрещены переходы между состояниями одинаковой четности (то есть между четными или между нечетными состояниями), а для дипольного магнитного и квадрупольного электрич. излучений - переходы между состояниями различной четности (то есть между четными и нечетными состояниями).



Наряду с точными О. п. по-J и $m_{J}$ существенны приближенные О. п. при дипольном излучении атомов для квантовых чисел, определяющих величины орбитальных и спиновых моментов электронов и проекций этих моментов. Напр., для атома с одним внешним электроном орбитальное (азимутальное) квантовое число $l$, определяющее величину орбитального момента электрона $\boldsymbol{M}_{l}\left[\boldsymbol{M}_{l}^{2}=\hbar^{2} l(l+1)\right]$, может изменяться на \textbackslash pm 1 ( $\Delta l=0$ невозможно, так как состояния с одинаковыми $l$ имеют одинаковую четность: они четные при четном $l$ и нечетные при нечетном $l$ ). Для сложных атомов квантовое число $L$, определяющее полный орбитальный момент всех электронов, подчинено приближенному О. п. $\Delta L=0, \pm 1$, а квантовое число $S$, определяющее полный спиновый момент всех электронов,- приближенному О. П. $\Delta S=0$, справедливому, если не учитывать спинорбитальное взаимодействие. Учет этого взаимодействия нарушает последнее О. п., и появляются так наз. интеркомбинационные переходы, вероятности к-рых тем больше, чем больше атомный номер элемента.





\end{document}